\section{Discussion and implications}

The central finding is that malicious-source success and misinformation exposure are related but distinct outcomes. A source can gain selection share and unique reposters by resembling traditional media, yet the network-wide false-repost share may barely move when users rapidly discount old information and social recommendations react negatively to false publications. Evaluations of platform policy should therefore report both source-level and content-level indicators. Removing one source metric from a dashboard can conceal displacement or exposure effects elsewhere.

Combined imitation produced the clearest audience gain and a delayed increase in false reposts. This suggests that policies based on a single surface characteristic are vulnerable to multidimensional camouflage. Source assessment should combine provenance, behavioural history, publication-level accuracy and coordination signals instead of relying on appearance alone. The result is compatible with evidence that source-quality judgments can support interventions \citep{pennycook2019}, while extending the question to adaptive source presentation.

Reach governance was the strongest directly controllable mechanism in the experiments. Lower direct visibility monotonically reduced the incremental effect of a malicious strategy, and zero reach eliminated it. This does not imply that account removal is always socially desirable or operationally feasible. Real sources can create new accounts, migrate across platforms, or exploit indirect sharing. The result instead establishes a controlled benchmark: when a platform can reduce initial recommendation and discovery exposure, malicious-source strategy has less opportunity to operate.

Social recommendation requires careful interpretation. Generic word of mouth could amplify harmful information, but the implemented mechanism is truth sensitive: false recommendations generate negative evidence and true recommendations positive evidence. Its protective effect is therefore a design result, not a description of X. A practical implication is that recommendation systems need reliable and timely integrity signals if they are expected to dampen misinformation. Delayed or noisy truth classification could yield different dynamics and should be examined in future scenarios.

Memory sensitivity places a boundary around policy claims. Long-lived evaluations amplify the downstream consequences of camouflage, while rapid decay prevents many effects from accumulating. Memory is not itself a policy variable in the same sense as reach, but interfaces, repeated exposure and ranking can influence how often source signals are refreshed. The sensitivity analysis is valuable precisely because it shows which conclusions survive behavioural uncertainty: audience capture is broadly robust; increased false exposure is not.

Several limitations remain. The network is random and small relative to X, users listen to a fixed number of contacts, and all agents use the same decision architecture. Fake-news probability is source-specific but content is not represented semantically. The Chilean survey supplies aggregate inputs without respondent-level uncertainty, detailed sampling documentation, or an external diffusion target. The platform's proprietary ranking, moderation and bot-detection systems are not reconstructed. Results are therefore mechanism-based counterfactuals rather than forecasts. Future work should incorporate empirical networks, user segments, delayed truth assessment, coordinated sources, account replacement, and calibration/validation against observed cascades while preserving the transparent source-behaviour distinction.
